% Generated by tools/render_uc_spec.py from usecase/系统规约.md.
% Do not treat local demo status as production approval.
\subsection{3.2 系统用例规约}
\noindent 下列十五项规约的语义以现行《系统规约》为准；旧 UC01～UC24 只按映射表迁移。
\subsubsection{3.2.1 UC-01 建立与查询顾客档案}\label{uc:01}
\begin{BSRUsecase}{UC-01 建立与查询顾客档案}
执行者与利益 & 顾客能管理自己的联系方式和尺码；门店租赁专员能识别回头客并避免重复录入；经理需要可审计的客户归属。 \\
关联 & FR-C01、FR-F01、DR-01、NFR-01。 \\
状态 & 演示已实现。 \\
触发 & 顾客注册/编辑资料，或专员协助现场开单时检索顾客。 \\
前置 & 注册时邮箱未被占用；检索时专员已登录且有门店业务权限。 \\
成功后置 & 生成或更新唯一 customer\_id；旧合同仍保留下单时快照。 \\
失败后置 & 不创建重复档案，不泄露匹配到的其他顾客资料。 \\
基本流程 & 1 顾客提交姓名、邮箱、电话、地址与可选量体资料。2 系统验证必填、格式和邮箱唯一性。3 系统生成不可复用的客户 ID，保存档案并记录时间。4 顾客可编辑本人资料；专员可用完整邮箱或客户 ID 查回头客并核对身份。5 新合同从当前档案生成独立快照。 \\
扩展 & 2a 邮箱已用：提示登录或找回，不显示该账号隐私。2b 电话/量体非法：指出字段，保留未提交表单。4a 搜索无结果：仅显示“未找到”；允许按授权流程新建。4b 专员无权限：拒绝并记审计。5a 顾客此后改资料：新资料用于未来订单，旧合同快照不变。 \\
字段/规则 & 邮箱用于精确匹配而非主键；姓名、电话、地址、出生年份、身高、体重、鞋码按当前需求可为空与否分别说明；每次租赁的技能自评属于 UC-04，不存为永久顾客等级。 \\
质量/约束 & 顾客仅见本人；敏感量体最小化展示；检索不能返回模糊匹配的完整隐私。 \\
\end{BSRUsecase}

\subsubsection{3.2.2 UC-02 按日期段查询可租装备}\label{uc:02}
\begin{BSRUsecase}{UC-02 按日期段查询可租装备}
执行者与利益 & 顾客先确定旅程时段，知道该时段还能租多少合适装备；专员可用同一库存口径协助。 \\
关联 & FR-C02/C03、FR-I01、BR-01/07、DR-02/03。 \\
状态 & 演示已实现，营业时间待 BSR 确认。 \\
触发 & 选择门店、起租日和归还日。 \\
前置 & 门店存在、开放预约；日期合法。 \\
成功后置 & 返回每规格在整个所选时段的可订量、适用尺码/款式与计价基础，不产生占用。 \\
失败后置 & 不显示误导性旧库存。 \\
基本流程 & 1 顾客选择门店及起止日。2 系统按挪威时区构成 09:00 至归还日 16:00 的租期并验证时长。3 顾客选择类别（板、鞋、雪杖、护膝、雪镜等）、子类、款式、尺码及“仅有货”。4 系统汇总符合规格且在整个区间可交付的实物，扣除重叠暂留/订单和未归还占用。5 展示剩余件数、租金和适用提示。 \\
扩展 & 2a 日期缺失、倒置、过长或已过：阻止查询并说明可改的字段。3a 配件无技能等级：不强迫选择技能。4a 某规格无货：显示零和可调整门店/日期，不自动替换规格。4b 结果查询后被他人占用：UC-04 再检验；旧数量失效。 \\
字段/规则 & store\_id、rental\_from/to、variant\_id、available\_qty、skill\_fit；可订量不是总库存。 \\
质量/约束 & 手机端可先选日期再筛选；空、加载、失败状态清晰；同一可用量算法供顾客和店员使用。 \\
\end{BSRUsecase}

\subsubsection{3.2.3 UC-03 管理购物车与透明报价}\label{uc:03}
\begin{BSRUsecase}{UC-03 管理购物车与透明报价}
执行者与利益 & 顾客可比较多件装备，清楚租金、押金、取还时刻和规则。 \\
关联 & FR-C04/C05、BR-02、DR-04/06。 \\
状态 & 演示已实现，正式价目待批准。 \\
触发 & 将规格加入购物车或修改数量。 \\
前置 & 已选有效门店与租期；规格可租。 \\
成功后置 & 购物车各行及 NOK 报价一致；未提交前不构成最终库存承诺。 \\
失败后置 & 原购物车草稿保留，不能显示过期报价为可提交。 \\
基本流程 & 1 加入规格与数量。2 系统展示类别/尺码、日租价、天数、行租金、按参考价值 8\% 算出的逐件押金。3 汇总租金、押金、其他已明确费用及应付金额，列取货与归还截止。4 顾客调整数量或移除行；系统重算。5 顾客进入规则确认。 \\
扩展 & 1a 数量超过当时可订量：拒绝并给出剩余数。2a 参考价值/价格缺失：不生成可付款报价，交经理维护价目。3a 定位或汇率不可用：仅显示 NOK；参考币种必须标汇率来源、时间与非合同效力。4a 改门店或租期：所有行重新询价和核量，原报价标失效。4b 购物车为空：不能结账。 \\
字段/规则 & qty 为正整数，金额为整数 NOK；历史订单锁定当时单价和政策，不随新价格变化。 \\
质量/约束 & 分项金额及总额同屏可核对；定位不能改权威账本。 \\
\end{BSRUsecase}

\subsubsection{3.2.4 UC-04 阅读规则并提交预约}\label{uc:04}
\begin{BSRUsecase}{UC-04 阅读规则并提交预约}
执行者与利益 & 顾客在知悉超时、损坏、押金和补差机制后预约；门店获得可追溯同意和短时库存承诺。 \\
关联 & FR-C06/C07、BR-01\textasciitilde{}05、DR-04/06。 \\
状态 & 演示已实现，政策及正式合同顺序待确认。 \\
触发 & 顾客点击继续结账，或门店专员为指定顾客现场开单。 \\
前置 & 实际提交者已登录；购物车/现场明细、完整档案、本次滑雪水平、租期和规则同意有效；创建请求带 16～128 字符的 Idempotency-Key。 \\
成功后置 & 保存规则文本版本、主动同意时间、报价和有限期暂留，生成订单号；实际提交账号、渠道和请求标识只映射一单。 \\
失败后置 & 不占用库存，不保存虚假的同意或请求映射。 \\
基本流程 & 1 展示损坏五级、正常磨损免责、超时计费、押金抵扣及超押金联系补差。2 顾客滚动浏览并主动勾选同意；现场由专员记录已向顾客出示并获同意。3 客户端为一次提交意向生成标识；服务端按实际提交账号和线上/现场渠道查询原映射，并核对稳定负载摘要。4 若没有原映射，系统记录规则版本、本次技能自评与报价快照，再次校验门店、租期、数量与库存。5 在同一事务建立暂留订单、请求映射、到期时刻及取货流程提示。 \\
扩展 & 2a 未滚动或未勾选：停留规则页。3a 缺失或非法提交标识：拒绝创建，提示重新提交；相同标识但内容改变：409 冲突，不能静默新建。3b 相同账号、渠道、标识和内容重试：200 返回原订单 ID、原报价及当前生效状态；暂留到期返回 EXPIRED，已取消返回 CANCELLED，均不延长暂留、不重复同意、通知或审计。3c 页面收到不确定网络错误或刷新：在同一登录会话内用原标识与原负载找回；若原单到期，保留购物车并引导重新选期，新意向须用新标识。4a 规则发布新版本：新意向重显最新文本并重新同意。4b 库存不足或价格变化：中止提交、回购物车并显示差异；不自动加价。5a 现场专员换目标顾客却复用标识：409；不同员工即使标识相同也按账号隔离。5b 同意后取消：保留审计，不视为已付款租赁。 \\
字段/规则 & reservation\_requests 保存实际 actor\_id、channel、request\_key、稳定请求摘要、reservation\_id 和服务器创建时刻；复合主键 (actor\_id,channel,request\_key)，一笔预约最多关联一条映射。标识由客户端随机生成，有效字符为 ASCII 字母、数字、点、下划线、冒号和连字符。规则版本、同意时间和顾客身份分开记录；浏览至末尾不等于已理解，主动勾选是必要但不替代正式签署。 \\
质量/约束 & 关键信息可在移动端完整阅读；禁止默认预勾选；同一标识并发请求只建一单，失败事务不占用标识；历史预约不伪造映射。当前单机 SQLite 事务已验证，生产多实例需单独验收。 \\
到期状态与后置条件 & 原始 HOLD 行在服务器时间达到 hold\_expires\_at 后，列表、详情、流程进度和同键重放统一返回生效状态 EXPIRED，但读取不改写原始行、金额、同意与合同快照。该订单不再占库存或进入员工未解决提醒；页面显示“暂留已过期”和历史 NOK 报价，不提供签约、付款或自行取消。重新预约可带回原装备、数量和门店；若购物车已有商品须先确认替换。顾客必须重新确认租期、核价和同意，并使用新的提交标识。服务端对旧订单的签约、付款、取消均返回冲突，不产生新交易。 \\
\end{BSRUsecase}

\subsubsection{3.2.5 UC-05 开具、变更并签署合同}\label{uc:05}
\begin{BSRUsecase}{UC-05 开具、变更并签署合同}
执行者与利益 & 顾客获得可打印、可复核的合同；专员保存承诺内容和签署证据；门店防止旧版被覆盖。 \\
关联 & FR-F02/F03/F06、BR-08、DR-04/06。 \\
状态 & 顾客自助确认与纸质见证的本地签付顺序已实现；合格线上签署为生产缺口。 \\
触发 & 顾客暂留预约后自助生成合同、员工现场开单，或租期内更换装备。 \\
前置 & 顾客和门店经办责任人可确定、暂留未过期、报价有效；签署对象为当前版本。 \\
成功后置 & 唯一合同号、版本和不可变条款快照；分别记录实际发起人与经办责任人及有效版本签署结果。 \\
失败后置 & 旧版保持可审计，未签版本不能演示付款或交付；拒绝请求不产生支付流水。 \\
基本流程 & 1 顾客在有效暂留订单中点击生成合同，或专员现场开单；系统把顾客、实际发起人、门店经办责任人、门店、日期、逐项规格/数量/单价/租金/押金/规则写入合同快照。2 系统分配唯一合同号及当前版本。3 顾客核对 NOK 金额、取还时间和规则。4 线上由顾客本人记录电子确认，现场由专员记录纸质签署见证；两者均为课程演示。5 系统保存签署方式、时间及快照摘要，打印视图显示来源和版本；随后才可进入 UC-06。 \\
扩展 & 1a 经办人未分配或关键字段缺失：阻止开具并提示联系门店。1b 顾客跨账号操作：拒绝。2a 重复请求：返回冲突并提示查看当前版本，不新建合同。3a 顾客要求更换：生成新版本并重新核价与签署，旧版只读。4a 拒签、错误快照摘要或暂留过期：不得演示付款和交付。4b 电子确认须明确标为演示，不声称法律有效签章。5a 旧版本回调：不生效并留审计。 \\
字段/规则 & 合同、订单、顾客、经办员工及政策均有恒定引用；contracts.created\_by 是实际发起人，快照中的 employee 是经办责任人，channel 区分线上自助与现场人工；顾客自助不计入 reservations.created\_by 员工开单业绩。 \\
质量/约束 & 后端签付守卫必须与页面一致；正式第三方支付和合格签署仍为外部集成目标。 \\
\end{BSRUsecase}

\subsubsection{3.2.6 UC-06 确认付款并生成备货任务}\label{uc:06}
\begin{BSRUsecase}{UC-06 确认付款并生成备货任务}
执行者与利益 & 顾客知道是否付款成功；门店只为有效订单备货。 \\
关联 & FR-C07、FR-F03/F04、DR-06、NFR-03。 \\
状态 & 演示付款已实现；真实支付为生产缺口。 \\
触发 & 当前合同签署后发起本地演示付款；正式流程将由第三方服务承接。 \\
前置 & 当前合同版本已签、暂留有效、订单未支付，首租宽限资格与签署快照一致。 \\
成功后置 & 登记一组租金与押金演示流水，仅成功后将订单转为备货队列；重复付款不追加流水或任务。 \\
失败后置 & 订单停留待付款/待人工核实，签署缺失时无流水，不显示“银行已到账”。 \\
基本流程 & 1 系统给出 NOK 应付总额、合同当前版和演示标识。2 顾客确认演示付款；服务端先核当前签署、暂留、首租宽限资格和未付状态。3 本地适配器分别写入租金与押金演示流水，不发生资金划转；正式目标改由独立第三方返回可验证回执。4 系统核对订单、金额、版本和请求引用。5 成功后记录支付状态并创建一份备货任务，通知有权限的门店岗位。 \\
扩展 & 2a 顾客取消：不生成支付成功记录。2b 未签、暂留过期或经办人缺失：拒绝付款，顾客重新生成有效预约或联系门店。2c 首租资格在暂留期间变化：要求重新预约并确认新合同。3a 正式目标中超时/回调缺失：标“待核实”，不重复扣款；由门店按对账流程处理。4a 金额或版本不符：拒绝入账、记录异常。5a 重复提交或回调：仅保留首笔有效租金及押金流水和一份备货任务。5b 本地演示：明确展示“模拟付款，无真实资金流”。 \\
字段/规则 & 交易引用、订单/合同版本、金额、币种、状态、回执时刻、演示/真实标识。 \\
质量/约束 & 外部失败可恢复；真实网关密钥与回调验证属于后续设计，不存入公开仓库。 \\
\end{BSRUsecase}

\subsubsection{3.2.7 UC-07 准备、核码并逐件交付}\label{uc:07}
\begin{BSRUsecase}{UC-07 准备、核码并逐件交付}
执行者与利益 & 专员知道拣货和验货责任；顾客可见每步进度，到店快速取走安全装备。 \\
关联 & FR-F04/F05、FR-C08、BR-07、DR-05。 \\
状态 & 演示已实现。 \\
触发 & 已确认订单进入门店备货队列。 \\
前置 & 订单有效，备货任务存在，店员获授权。 \\
成功后置 & 指定实物与订单行绑定、完成检查并逐件记录交付人/时刻；顾客看到状态。 \\
失败后置 & 不可交付实物不转为已借出。 \\
基本流程 & 1 经理或授权专员指派备货人。2 专员按订单规格逐件扫描/选择库存号。3 核对尺码、清洁、外观和适用安全项，记录检查。4 全部可交付时标“已备妥”，向顾客同步具名进度。5 顾客到店提供取货码，专员核身份、合同当前版签署和实际装备。6 系统逐件记录交付时刻、处理人，状态转“租赁中”。 \\
扩展 & 2a 库存已占用/跨店/隔离：拒绝绑定，标阻塞并通知负责人。3a 检查不合格：隔离实物，重新拣货；原记录留痕。4a 只完成部分：显示准备进度，不谎称全部可取。5a 取货码错/失效/已用：拒绝交付，走人工核实；不能仅凭码绕过身份核对。5b 未签当前合同：返回 UC-05。6a 重复扫描：不产生二次借出。 \\
字段/规则 & 每件库存号、检查项、状态、处理人、事件时间、顾客可见说明和内部备注分开。 \\
质量/约束 & 授权角色间同步；顾客看具名负责人和进度，不能看员工内部隐私。 \\
\end{BSRUsecase}

\subsubsection{3.2.8 UC-08 租期内更换装备与部分归还}\label{uc:08}
\begin{BSRUsecase}{UC-08 租期内更换装备与部分归还}
执行者与利益 & 顾客可换装或先归还部分装备；门店维护单件与差价准确性。 \\
关联 & FR-F06/F07、FR-D04、DR-05/06。 \\
状态 & 演示已实现。 \\
触发 & 顾客提出升级、降级、换尺码或分批归还。 \\
前置 & 合同有效且至少一件已借出；员工获授权；拟换单件可租。 \\
成功后置 & 仅目标单件的交接和费用发生变化；未变更单件继续有效，合同产生新版本。 \\
失败后置 & 不出现旧件已收回而新件未交付却被记成完整换装。 \\
基本流程 & 1 专员选定原单件及调整原因。2 系统预览剩余租期、旧/新价格及差额，展示应收或应退。3 顾客确认。4 系统检查当前合同版本和新件可用性。5 逐件记录旧件实还、检查，新件安全检查与交付。6 生成合同新版本及模拟差额流水，保留原版。 \\
扩展 & 2a 价格政策缺失：停止报价，经理维护。3a 顾客拒绝：原合同不变。4a 并发占用或旧版本：重算并要求重新确认。5a 旧件损坏：建立鉴定，不把疑似费用直接并入换装差价。6a 仅部分归还：已归还件冻结实还时刻与超时费；未归还件仍占库存、押金不结清。6b 重复提交：只生成一次调整版本。 \\
字段/规则 & 调整原因、原/新 unit\_id、差额 NOK、顾客确认、合同版本及逐件交接时间。 \\
质量/约束 & 预览不等于执行；演示资金流明示非真实收退。 \\
\end{BSRUsecase}

\subsubsection{3.2.9 UC-09 逐件归还与跨店接收}\label{uc:09}
\begin{BSRUsecase}{UC-09 逐件归还与跨店接收}
执行者与利益 & 顾客知道归还已被确认；专员掌握实物去向和异常；经理能核对跨店资产。 \\
关联 & FR-F07/F08、FR-I02、BR-07、DR-03/05。 \\
状态 & 演示已实现。 \\
触发 & 顾客将一件或多件装备交到原店或异店。 \\
前置 & 单件属于 ACTIVE 租赁且尚未归还；接收员工有当前所属店权限；异店须输入完整订单号或取货码。 \\
成功后置 & 每件记录实还时间、员工所属接收门店、处理人和验收事实；仅相关件释放借出状态；异店员工只收到本批受理回执。 \\
失败后置 & 未核实实物继续标“未归还”，不能提前退款；无业务写入或通知，也不暴露外店逐单资料。 \\
基本流程 & 1 本店专员查本店订单；异店专员凭完整且正确的代码精确核对在租实物。2 系统对本店展示授权订单资料，对异店仅展示归还截止和未还单件编号、名称、尺码及当前接收店。3 专员逐件记录“完好/需清洁/疑似损坏”事实及备注；服务端重验代码、未还单件和当前账号所属店。4 系统以服务器时刻记录实还，冻结该件超时金额并更新实物当前所在门店。5 完好件经清洁检查恢复可租；疑似损坏件进入 UC-10；仍有未归还件则订单维持进行中。6 系统给异店员工返回不含顾客、押金、合同和费用的本批接收回执，原店继续负责结算。 \\
扩展 & 1a 代码不完整、错误、订单非 ACTIVE 或外店按 ID 直接访问：拒绝，失败状态/文案不区分外店对象存在性。2a 实物不属于订单或已归还：拒绝重复入库或重复计费，原接收记录仅原店/本人可查。3a 员工伪造 returnStoreId 为其他门店或提交额外单件：拒绝且不写入。3b 店员不能确定损伤：标待鉴定，不直接收费。4a 异店归还只更新当前门店，不自动改变永久归属；归属转移走 UC-13。5a 部分归还：顾客可见逐件进度与尚欠件。5b 逾期未还：持续占用未来可订量并进入原店负责人提醒。 \\
字段/规则 & returned\_at、return\_store\_id、inspection\_condition、receiver\_id、late\_fee\_nok；异店查找只返订单标识/编号、状态、截止和未还单件，提交重验精确码。 \\
质量/约束 & 跨角色读同一事实但按最小授权投影；归还与状态写入应原子完成，拒绝无副作用；异店回执不含单笔财务及私有证据。 \\
\end{BSRUsecase}

\subsubsection{3.2.10 UC-10 受损初评、照片鉴定与经理复核}\label{uc:10}
\begin{BSRUsecase}{UC-10 受损初评、照片鉴定与经理复核}
执行者与利益 & 顾客获得可核对的损伤证据和计价依据；技师明确任务；经理控制误收费。 \\
关联 & FR-D01\textasciitilde{}D03、BR-04/06、DR-07。 \\
状态 & 演示已实现，商业收费与争议程序待批准。 \\
触发 & UC-09 标记疑似损坏。 \\
前置 & 有归还事实、单件已隔离、鉴定任务已指派。 \\
成功后置 & 24 小时内记录初评反馈；经理核定或退回，有版本化证据；核定前不生成确定损坏费。 \\
失败后置 & 超时或缺证据持续隔离，提醒经理与顾客，不自动判定顾客有责。 \\
基本流程 & 1 系统从发现时间计算初评截止并将任务分配给一位技师。2 技师核实库存号、受损部位，上传对应照片。3 技师描述事实，标无损/正常磨损、轻微、中度、严重或不可修复，填写工时、材料、折旧估值依据和初评。4 系统给顾客显示可公开证据和进度。5 经理审阅实物、照片与计价依据，核定金额或退回补证。6 系统保留历史版本、复核人和时间。 \\
扩展 & 1a 技师无法到岗：经理改派，原指派留痕。2a 无照片或照片与单件不符：不能提交收费鉴定；不得含面部/证件等无关资料。3a 正常磨损：费用为零。3b 不可修复：以有依据的折旧后重置估值计，不按等级乘固定罚款比例。5a 经理认为证据不足：退回补充，旧照片/意见不覆盖。5b 顾客提出异议：标争议并暂停 UC-11 最终结清。5c 超过发现后24小时：提示逾期并留原因，24小时目标仅指初评反馈，不指维修完成。 \\
字段/规则 & 单件、部位、照片、级别、费用构成、估值依据、指派/反馈/复核时间和状态。 \\
质量/约束 & 照片只允许订单本人及相关授权员工/技师访问；提取原图定位信息应避免隐私泄露。 \\
\end{BSRUsecase}

\subsubsection{3.2.11 UC-11 结算超时、损坏与押金}\label{uc:11}
\begin{BSRUsecase}{UC-11 结算超时、损坏与押金}
执行者与利益 & 顾客看清租金外费用、押金扣退和补差；经理防误扣；门店闭环账务。 \\
关联 & FR-D04/D05、FR-F08、BR-03\textasciitilde{}05、DR-06。 \\
状态 & 演示结算已实现，费率/退款时限/真实到账待确认。 \\
触发 & 全部实物归还且异常鉴定有结果。 \\
前置 & 单件实还时间齐全，所有损坏已核定或标争议，订单适用规则版本可查。 \\
成功后置 & 分项结算、押金抵扣、可退额或补差额冻结；无争议且补差闭环后发起退款。 \\
失败后置 & 资金状态不假报成功，争议/支付失败不标完成。 \\
基本流程 & 1 系统按每件实还时间和订单冻结的超时规则计算超时费。2 经理确认核定损坏费与依据。3 系统分列各件超时费、损坏费及总费用。4 计算 抵扣=min(总费用,押金)、可退=max(押金-总费用,0)、补差=max(总费用-押金,0)。5 顾客可查明细与证据。6 零补差、无争议时发起退款；补差时专员联系顾客、记录结果与分次凭证，再核对结清。 \\
扩展 & 1a 首笔已生效租赁：适用订单冻结的120分钟宽限；仅暂留或取消单不消耗资格。1b 旧单无超时规则版本：不追溯自动收费，交人工核实。2a 损坏未核定：暂停最终结算。4a 总费用大于押金：不得自动从其他渠道扣款。5a 顾客异议：保留证据，暂停最终退款和补差结清。6a 退款失败/回执不明：标失败或待核实并由经理处理，不标已到账。6b 分次补差重复凭证或超额：拒绝并留审计。 \\
字段/规则 & 订单级押金、单件超时、单件损坏、退款/补差流水、幂等引用、争议状态。 \\
质量/约束 & 整数 NOK、逐项透明、历史政策不回写；演示资金状态与真实银行状态分离。 \\
\end{BSRUsecase}

\subsubsection{3.2.12 UC-12 发起维修、技师处理与安全验收}\label{uc:12}
\begin{BSRUsecase}{UC-12 发起维修、技师处理与安全验收}
执行者与利益 & 门店移走不安全装备；技师按任务记录劳动与材料；经理控制再出租安全。 \\
关联 & FR-R01/R02、FR-D03、BR-06/07、DR-07。 \\
状态 & 演示已实现，外部通知渠道待选。 \\
触发 & 设备故障、受损鉴定或定期维护。 \\
前置 & 实物存在且当前不在可交付占用中；发起人有权限。 \\
成功后置 & 一单一件一责任技师；维修记录和费率快照齐全，经理安全验收通过后才恢复可租。 \\
失败后置 & 设备继续隔离，不会因技师点击完成而自动上架。 \\
基本流程 & 1 专员或经理描述故障、选单件与技师。2 系统创建工单、冻结技师当时小时费率并发送应用内任务提醒。3 技师诊断，登记工时（0.25小时倍数）、配件、处理和备注。4 技师完成适用安全检查，提交经理验收。5 经理验收通过或退回重修；通过后再确认清洁并恢复可租。 \\
扩展 & 1a 设备已在另一未完工单或被借出：拒绝冲突工单，先核对实物。2a 技师改派：保存原指派和新费率/责任的明确生效规则，不改写已发生工时。3a 需等配件：标等待配件并持续隔离。3b 工时不是15分钟倍数：拒绝提交。4a 安全项缺失：不能申请验收。5a 经理退回：保留本轮诊断与检查，继续隔离。5b 判定报废：状态为报废，可订量保持零。 \\
字段/规则 & 工单号、单件、责任技师、报修人、故障、工时、材料、费率版本、诊断、验收结果。 \\
质量/约束 & 技师只见本人任务和必需客户背景；外部消息实际送达不可凭应用内提醒宣称。 \\
\end{BSRUsecase}

\subsubsection{3.2.13 UC-13 新增门店装备、采购入库与跨店调拨}\label{uc:13}
\begin{BSRUsecase}{UC-13 新增门店装备、采购入库与跨店调拨}
执行者与利益 & 经理保持多店资产位置可信；顾客查询到正确门店的真实可租量。 \\
关联 & FR-I01/I02、DR-02/03、BR-07。 \\
状态 & 演示已实现，真实多实例容量待验收。 \\
触发 & 新增营业点/采购装备，或将装备永久调往另一店。 \\
前置 & 经理获授权；目标店存在且设备未被借出、鉴定或维修。 \\
成功后置 & 实物编号、规格、成本及门店信息完整；调拨发送/接收流水一致。 \\
失败后置 & 不能出现两店同时认为同一件归属自己。 \\
基本流程 & 1 授权管理岗位维护门店主数据和类别/规格。2 经理仅向所属店采购入库，逐件分配唯一库存号、购入日/成本/参考价值及固定适用等级。3 经理仅从所属店发起调拨，设备转为在途并退出可订量。4 目标店所属员工核对库存号并确认接收。5 系统更新当前门店和永久归属，保存发送/接收人及时刻。 \\
扩展 & 2a 编号重复、必填缺失或拟入库店非经理所属店：拒绝入库。2b 护膝、雪镜等配件不强制滑雪技能等级。3a 已占用/维修中或来源店非经理所属店：拒绝调拨。4a 实物不符/未收到或当前员工不属于接收店：保持在途，不能代收。5a 异店归还但非永久调拨：只改变当前门店，不执行本用例归属变更。 \\
字段/规则 & 门店、类别父子关系、规格、单件、采购值、当前店、归属店、调拨事件。 \\
质量/约束 & 可新增门店，不硬编码两店；库存跨角色读取一致。 \\
\end{BSRUsecase}

\subsubsection{3.2.14 UC-14 排班、上下班打卡与年度业绩}\label{uc:14}
\begin{BSRUsecase}{UC-14 排班、上下班打卡与年度业绩}
执行者与利益 & 员工知道今日岗位；经理掌握计划和实际到岗；经营负责人看到年度开单归属。 \\
关联 & FR-W01\textasciitilde{}W03、DR-08。 \\
状态 & 演示排班/打卡/报表已实现；奖金口径待 BSR 确认。 \\
触发 & 经理排班、非顾客员工打卡或负责人查询报表。 \\
前置 & 员工档案有效；经理有门店授权；年度报表只读权限已验证。 \\
成功后置 & 计划班次、实际到岗和合同归属各自独立、可查询。 \\
失败后置 & 不能用排班记录伪造打卡，也不自动发放奖金。 \\
基本流程 & 1 经理按日期、门店、员工设置 08:00–22:00 内一组起止时刻。2 系统检查该员工当日最多一班及冲突。3 员工本人上下班打卡，服务器记录挪威当地对应时刻。4 经理每日查看计划、到岗、迟到和早退。5 经营负责人按挪威年度查看已确认合同创建人的净租金归属与汇总。 \\
扩展 & 1a 经理安排外店员工或修改外店班次：拒绝，跨店支援另需明确授权。2a 已有冲突/已打卡班次需改：拒绝直接覆盖，走经理核实更正并留痕。3a 重复上下班打卡：保留原事实，不重写时间。3b 缺卡：明确标缺卡待核实。5a 合同取消或退款：报表按统一净额口径重算并标规则版本。5b 员工跨店支援：显示所属店与当日工作店，但当前演示不自动授予外店操作权。5c 奖金比例未批准：只报告依据，不生成应发工资。 \\
字段/规则 & 员工 ID、岗位/门店、班次日期、计划起止、实际打卡、合同创建人、年度和净租金。 \\
质量/约束 & 个人隐私按岗位展示；淡旺季 2–15 人容量尚需真实规模验证。 \\
\end{BSRUsecase}

\subsubsection{3.2.15 UC-15 新增/未解决提醒、权限与审计}\label{uc:15}
\begin{BSRUsecase}{UC-15 新增/未解决提醒、权限与审计}
执行者与利益 & 非顾客岗位看见自己要处理的新增与未结事项；管理员可追溯越权和关键变更。 \\
关联 & FR-O01/O02、FR-A01、NFR-01/03。 \\
状态 & 演示已实现。 \\
触发 & 新订单、逾期、维修指派、鉴定退回、补差争议、责任人变更或用户登录。 \\
前置 & 事件已保存；查看者已认证。 \\
成功后置 & 各角色只见授权事项，新增标记与业务未解决状态独立；审计保留写入历史。 \\
失败后置 & 无权限时不泄漏订单摘要、顾客身份或证据。 \\
基本流程 & 1 系统把事件归于原订单门店、指派技师或经营汇总角色。2 非顾客界面相应位置显示小红色感叹号和待办数量。3 店员/经理只见本店逐笔提醒并可打开有权详情；技师只见本人任务；经营负责人/管理员只见按门店、状态聚合数量，可标记本人汇总已阅。4 已读只清除本人“新增待阅”，业务未闭环时仍列“未解决”。5 改派、经理退回或新事件重新提醒新负责人。6 管理员仅查看账号、员工及门店配置的安全审计元数据；经理业务审计限本店。 \\
扩展 & 1a 技师未受派：不见顾客订单。1b 外店员工按订单 ID 直接请求提醒/已读或详情：与不存在对象统一拒绝，不改变已读状态。2a 实时连接中断：回到页面重新拉取，不能凭旧缓存宣布完成。3a 用户已读但未处理：未解决标记保留。3b 经营/管理员请求逐单提醒或逐单已读：拒绝，聚合只含数量。4a 完成/取消/过期：按实际状态退出待办，历史审计仍在。5a 原负责人离职：经理重派并保留过程。6a 管理员请求代做经理复核：除非另有业务授权，否则拒绝。 \\
字段/规则 & 事件对象、原订单门店、版本、责任人、每账号阅读状态、业务关闭状态、汇总门店/状态和审计操作者/时刻。 \\
质量/约束 & 敏感正文不放跨店通知或汇总响应；业务动作必须重新核对象归属；编辑草稿不被异步更新覆盖。 \\
\end{BSRUsecase}

